\documentclass[conference]{IEEEtran}
\IEEEoverridecommandlockouts
\usepackage{tabularx}
\usepackage{booktabs}
\usepackage{cite}
\usepackage{amsmath,amssymb,amsfonts}
\usepackage{algorithmic}
\usepackage{graphicx}
\usepackage{textcomp}
\usepackage{xcolor}
\usepackage{newtxtext,newtxmath}

\def\BibTeX{{\rm B\kern-.05em{\sc i\kern-.025em b}\kern-.08em
    T\kern-.1667em\lower.7ex\hbox{E}\kern-.125emX}}

\begin{document}

\title{Comparing Critic and User Movie Reviews Using Data Mining}


\author{\IEEEauthorblockN{1\textsuperscript{st} Brady Lauritsen}
  \IEEEauthorblockA{\textit{University of Nebraska-Lincoln} \\
    \textit{Computer Science}\\
    Lincoln, NE, USA}
  \and
  \IEEEauthorblockN{2\textsuperscript{nd} Sagun Karki}
  \IEEEauthorblockA{\textit{University of Nebraska-Lincoln} \\
    \textit{Computer Science \& Data Science}\\
    Lincoln, NE, USA}
  \and
  \IEEEauthorblockN{3\textsuperscript{rd} Harrison Johs}
  \IEEEauthorblockA{\textit{University of Nebraska-Lincoln} \\
    \textit{Computer Science}\\
    Lincoln, NE, USA}
  \and
  \IEEEauthorblockN{4\textsuperscript{th} Max Buettenback}
  \IEEEauthorblockA{\textit{University of Nebraska-Lincoln} \\
    \textit{Computer Science}\\
    Lincoln, NE, USA}
}

\maketitle

\begin{abstract}
  Movie ratings from professional critics and general audiences often differ significantly.
  This project explores the relationship between critic reviews from Rotten Tomatoes and
  user reviews from Letterboxd using data mining techniques. The goal is to identify patterns,
  analyze sentiment differences, and determine whether critic ratings can predict user ratings.
  By combining multiple datasets and performing preprocessing, correlation analysis, and
  predictive modeling, we aim to better understand how critics and audiences evaluate films.
\end{abstract}

\begin{IEEEkeywords}
  data mining, sentiment analysis, movie reviews, correlation analysis, prediction
\end{IEEEkeywords}

\section{Introduction}
The difference between professional critic reviews and general audience
reviews has long been a topic of discussion in the film industry. Critics
often evaluate films based on technical quality, artistic merit, and
narrative structure, while general audiences may prioritize entertainment
value and personal enjoyment. These differences can lead to noticeable
discrepancies in ratings.

This project aims to analyze and compare critic ratings from Rotten Tomatoes
with user ratings from Letterboxd using data mining techniques. By combining
these datasets, we seek to identify patterns of agreement and disagreement,
analyze sentiment differences in written reviews, and determine whether
critic scores can be used to predict user ratings.

\section{Objectives}

The goals of this project are as follows:

\begin{itemize}

  \item \textbf{Objective 1 -- Data Preprocessing:}
        Develop a unified dataset by cleaning and integrating Rotten Tomatoes
        critic review data and Letterboxd user rating data. This includes handling
        missing values, removing duplicates, normalizing rating scales, and matching
        movies by title and release year. Review text will be preprocessed using
        tokenization, lowercasing, and stop-word removal.

  \item \textbf{Objective 2 -- Correlation Analysis:}
        Measure statistical relationships between critic ratings and user ratings,
        identifying movies with strong agreement or disagreement and analyzing
        patterns across genres.

  \item \textbf{Objective 3 -- Sentiment Analysis:}
        Perform sentiment analysis on review text to compare the emotional tone of
        critic reviews and user reviews.

  \item \textbf{Objective 4 -- Predictive Modeling:}
        Develop machine learning models to predict user ratings from critic scores
        and movie attributes, and evaluate prediction performance.

\end{itemize}

\section{Datasets}

\subsection{Dataset 1: Rotten Tomatoes}
\begin{itemize}
  \item \textbf{Source:} Kaggle \cite{leone2020rotten}
  \item \textbf{Size:} 17,000+ movies and 1M+ critic reviews.
  \item \textbf{Scope:} Data scraped as of October 2020.
\end{itemize}

\begin{table}[htbp]
  \caption{Rotten Tomatoes Key Attributes}
  \begin{center}
    \begin{tabularx}{\columnwidth}{|l|l|X|}
      \hline
      \textbf{Attribute}  & \textbf{Data Type} & \textbf{Description}               \\ \hline
      movie\_title        & Nominal            & The official title of the film.    \\ \hline
      genres              & Nominal            & Categorical genre tags.            \\ \hline
      tomatometer\_rating & Ratio              & Professional critic score (0-100). \\ \hline
      review\_content     & Nominal            & Raw text of the critic review.     \\ \hline
      release\_year       & Interval           & The year the movie was released.   \\ \hline
    \end{tabularx}
    \label{tab1}
  \end{center}
\end{table}

\subsection{Dataset 2: Letterboxd}
\begin{itemize}
  \item \textbf{Source:} Kaggle \cite{letterboxd2020}
  \item \textbf{Size:} Approximately 2.5 million user ratings.
  \item \textbf{Scope:} User-submitted ratings on a 0.5--5 star scale.
\end{itemize}

\begin{table}[htbp]
  \caption{Letterboxd Key Attributes}
  \begin{center}
    \begin{tabularx}{\columnwidth}{|l|l|X|}
      \hline
      \textbf{Attribute} & \textbf{Data Type} & \textbf{Description}                 \\ \hline
      movie\_title       & Nominal            & The title used for dataset matching. \\ \hline
      rating\_val        & Ordinal            & User score from 0.5 to 5.0.          \\ \hline
      user\_id           & Nominal            & Unique identifier for the user.      \\ \hline
      year\_released     & Interval           & Release year for matching purposes.  \\ \hline
    \end{tabularx}
    \label{tab2}
  \end{center}
\end{table}

\section{Data Processing}

Before analysis, both datasets will undergo preprocessing to ensure
consistency and compatibility. This includes cleaning movie titles,
removing duplicates, and handling missing values. Rating scales will
be normalized to allow direct comparison between critic and user scores.

Movies will be matched across datasets using title and release year.
Text reviews will be preprocessed using tokenization, lowercasing,
removal of stop words, and punctuation filtering to prepare for
sentiment analysis. Feature selection will be performed to identify
relevant variables for predictive modeling.

\section{Methods}

To accomplish the project objectives, several data mining techniques
will be applied. Correlation analysis will be performed using Pearson
and Spearman correlation coefficients to measure relationships between
critic and user ratings.

Sentiment analysis will be conducted on review text using natural
language processing techniques, including polarity scoring and
lexicon-based methods.

For predictive modeling, machine learning algorithms such as linear
regression and random forest regression will be implemented to predict
user ratings based on critic scores and movie attributes. Model
performance will be evaluated using standard regression metrics.

\section{Evaluation}

Project performance will be evaluated using multiple quantitative
metrics. Correlation strength will be measured using Pearson and
Spearman coefficients. Predictive models will be evaluated using
Root Mean Squared Error (RMSE) and R-squared values.

Additionally, visualizations such as scatter plots and distribution
graphs will be used to interpret relationships between critic and
user ratings.

\section{Implementation Plan}

\subsection{Tools}
The project will be implemented using Python. Libraries such as
Pandas and NumPy will be used for data preprocessing, Scikit-learn
for machine learning, and Matplotlib or Seaborn for visualization.
Natural Language Toolkit (NLTK) or TextBlob will be used for
sentiment analysis.

\subsection{Timeline}
\begin{itemize}
  \item Week 1--2: Data collection and preprocessing
  \item Week 3: Dataset integration and exploratory analysis
  \item Week 4: Correlation and sentiment analysis
  \item Week 5: Predictive modeling
  \item Week 6: Evaluation, visualization, and final report preparation
\end{itemize}

\subsection{Meeting Schedule}

The team will meet virtually on a weekly basis on Wednesday in the afternoon(sometime between 2-4 pm)

\section{Team Contributions}

\begin{itemize}
  \item Brady Lauritsen: Responsible for predictive modeling implementation, including the development of linear and random forest regression models and performance evaluation.
  \item Sagun Karki: Responsible for sentiment analysis implementation, finding correlations among attributes, and creating reports and visualizations.
  \item Max Buettenback: Responsible for data preparation and cleaning and exploratory analysis.
  \item Harrison Johs: Responsible for data preparation and cleaning and exploratory analysis.
\end{itemize}

\section{Conclusion}

This project aims to provide a data-driven comparison between
professional critic ratings and general audience ratings. By
integrating multiple datasets and applying data mining techniques,
we seek to identify patterns of agreement and disagreement, analyze
sentiment differences, and evaluate predictive relationships between
critic and user evaluations. The results will contribute to a better
understanding of how different groups assess film quality.

\bibliographystyle{IEEEtran}
\bibliography{references}
\end{document}